\chapter{Wittig and Other C=C Forming Reactions}\label{ch:b2:wittig}

A female housefly attracts males with a hydrocarbon of 23 carbons carrying a single double bond, between
carbons 9 and 10, with the (\textit{Z}) geometry. To make it in quantity, a chemist must put that one
double bond exactly there, and with that geometry: an elimination would scatter it over several
positions and give both geometries. This chapter builds double bonds the other way, by joining two
pieces at a carbonyl group, so that the double bond appears where the \ce{C=O} was.
% ledger: use:muscalure

\begin{recall}
\Cref{ch:b2:enolates-aldol}: the \omterm{def:b2:enolates-aldol:aldol}{aldol condensation}. \Cref{ch:b2:alkene-redox}: the hydrogenation of
an alkyne to a (\textit{Z})-alkene over a poisoned catalyst. The Year 1 volume: E1, E2 and Zaitsev's
rule, bimolecular substitution, Grignard reagents and polarity inversion, the
(\textit{Z})/(\textit{E}) descriptors. The school volume: atom economy.
\end{recall}

\section{Phosphonium salts and ylides}

\begin{definition}[Phosphonium ylides]\label{def:b2:wittig:ylide}
A \emph{phosphonium salt}\index{phosphonium salt} $\mathrm{R_3P^+{-}CH_2R'\ X^-}$ is made by bimolecular
substitution of a halogenoalkane by a phosphine, usually triphenylphosphine \ce{Ph3P}. Removing a
hydrogen from the carbon bound to phosphorus gives a \emph{phosphonium ylide}\index{phosphonium ylide},
a neutral molecule with adjacent opposite charges, $\mathrm{Ph_3P^+{-}CH^-{-}R'}$, also written $\mathrm{Ph_3P{=}CHR'}$. A
\emph{stabilised ylide}\index{stabilised ylide} carries on that carbon an electron-withdrawing group
(an ester, a ketone, a \omterm{def:b2:amines:nitrile}{nitrile}) that delocalises the negative charge.
\end{definition}

\begin{omfigure}
\setchemfig{atom sep=1.9em}
\schemestart
\chemfig{Ph_3P} \+ \chemfig{CH_3-I}
\arrow{->[S$_\mathrm{N}$2]}[,0.9]
\chemfig{Ph_3P^{+}-CH_3}\;\chemfig{I^{-}}
\arrow{->[\ce{BuLi}][$-$ \ce{BuH}, \ce{LiI}]}[,1.3]
\chemfig{Ph_3P^{+}-\chemabove{C}{\scriptstyle -}H_2}
\arrow{<->}[,0.8]
\chemfig{Ph_3P=CH_2}
\schemestop
\par\vspace{6pt}

{\small Methylidenetriphenylphosphorane, the simplest ylide. Triphenylphosphine displaces iodide
from iodomethane; butyllithium removes a proton from the methyl group. The two structures on the right
are two ways of writing the same molecule: the charge-separated one is the closer to reality.}
\end{omfigure}

\begin{proposition}[Making the ylide]\label{prop:b2:wittig:ylide-acidity}
A \omterm{def:b2:wittig:ylide}{phosphonium salt} is far more acidic at the carbon next to phosphorus than an alkane. A simple alkyl
\omterm{def:b2:wittig:ylide}{phosphonium salt} still needs a very strong base (butyllithium, sodium hydride, sodium amide) under dry,
inert conditions; a salt whose \ce{CH2} also carries a carbonyl group is deprotonated by a weak base
(sodium carbonate, sodium hydroxide), and the \omterm{def:b2:wittig:ylide}{stabilised ylide} formed can be isolated and stored.
\end{proposition}

\begin{proof}[Argument]
The positive phosphorus next to the carbanion stabilises it by its charge and by the polarisability of
the large phosphorus atom: a simple ylide is a stabilised carbanion, but less so than an \omterm{def:b2:enolates-aldol:enolate}{enolate},
hence the strong base. An adjacent carbonyl spreads the negative charge onto oxygen, as in an \omterm{def:b2:enolates-aldol:enolate}{enolate}
(\cref{prop:b2:enolates-aldol:alpha-acidity}), on top of the phosphorus: the conjugate acid becomes
acidic enough for a weak base, and the ylide, much less basic, no longer reacts with water or air.
\end{proof}

\section{The Wittig reaction}

\begin{definition}[Wittig reaction]\label{def:b2:wittig:wittig}
The \emph{Wittig reaction}\index{Wittig reaction} of a \omterm{def:b2:wittig:ylide}{phosphonium ylide} with an aldehyde or a ketone
gives an alkene and triphenylphosphine oxide:
\begin{equation*}
\mathrm{Ph_3P{=}CHR + R'CHO \longrightarrow R'CH{=}CHR + Ph_3P{=}O}.
\end{equation*}
It goes through an \emph{oxaphosphetane}\index{oxaphosphetane}, a four-membered ring of phosphorus, two
carbons and oxygen, formed by the ylide carbon bonding to the carbonyl carbon while the carbonyl oxygen
bonds to phosphorus.
\end{definition}

\begin{omfigure}
\setchemfig{atom sep=1.9em}
\schemestart
\chemfig{Ph_3P=CH-R} \+ \chemfig{O=CH-R'}
\arrow{->[addition]}[,1.0]
\chemfig{Ph_3P?-[6]CH(-[4]R)-[0]CH(-[0]R')-[2]O?}
\schemestop

\medskip
\schemestart
\arrow{->[elimination]}[,1.1]
\chemfig{R-CH=CH-R'} \+ \chemfig{Ph_3P=O}
\schemestop
\par\vspace{6pt}

{\small The \omterm{def:b2:wittig:wittig}{Wittig reaction}. The ylide carbon and the carbonyl carbon bond together while oxygen bonds
to phosphorus, in one step: the \omterm{def:b2:wittig:wittig}{oxaphosphetane}. Its ring then opens the other way, breaking the
\ce{P-C} and \ce{C-O} bonds: the \ce{C=C} bond forms between the two former reacting carbons, and the
\ce{P=O} bond of triphenylphosphine oxide forms at the same time.}
\end{omfigure}

\begin{proposition}[A double bond exactly in place]\label{prop:b2:wittig:regiospecific}
In a \omterm{def:b2:wittig:wittig}{Wittig reaction} the new \ce{C=C} bond joins the former carbonyl carbon to the former ylide carbon,
and nowhere else.
\end{proposition}

\begin{proof}[Argument]
The only bonds made and broken are those of the four-membered ring: the ylide carbon to the carbonyl
carbon, then the \ce{C-P} and \ce{C-O} bonds. No hydrogen moves and no carbocation or carbanion forms
along the chain, so the double bond cannot shift. An elimination, by contrast, chooses between the
hydrogens on several neighbouring carbons (Zaitsev's rule), and an E1 reaction may rearrange.
\end{proof}

\begin{proposition}[Driving force]\label{prop:b2:wittig:driving}
The formation of the \ce{P=O} bond of triphenylphosphine oxide makes the \omterm{def:b2:wittig:wittig}{Wittig reaction} strongly
favourable and irreversible.
\end{proposition}

\begin{proof}[Argument]
The reaction trades a \ce{C=O} and a \ce{P-C} bond (in the ylide, the \ce{P=C} is mostly a
\ce{P+-C-} single bond) for a \ce{C=C} and a \ce{P=O} bond. The \ce{P=O} bond of a phosphine oxide is
very strong, much stronger than the \ce{P-C} bond lost, and the \ce{C=O} lost is compensated in part
by the \ce{C=C} formed: the balance is clearly downhill. Phosphorus is the oxygen sink of the
reaction.
\end{proof}

\section{Stereoselectivity}

\begin{proposition}[Geometry of the alkene]\label{prop:b2:wittig:stereo}
With an aldehyde, a non-stabilised ylide (\ce{Ph3P=CHR}, R an alkyl group) gives mainly the (\textit{Z})-alkene,
under salt-free conditions. A \omterm{def:b2:wittig:ylide}{stabilised ylide} (\ce{Ph3P=CH-COOR}, \ce{Ph3P=CH-COR}) gives mainly the
(\textit{E})-alkene. An ylide stabilised only by an aryl group gives mixtures.
\end{proposition}

\admitted

These are observations. Their usual explanation reads the \omterm{def:b2:wittig:wittig}{Wittig reaction} as a competition between
rates and stabilities. With a reactive, non-stabilised ylide, the \omterm{def:b2:wittig:wittig}{oxaphosphetane} forms fast and
irreversibly, through the less crowded transition state, in which the two substituents end up on the
same side of the ring; it then collapses with retention of that arrangement: kinetic control gives the
(\textit{Z})-alkene. With a \omterm{def:b2:wittig:ylide}{stabilised ylide}, the first step is slower and reversible; the
\omterm{def:b2:wittig:wittig}{oxaphosphetanes} interconvert through the starting materials, and the one with the substituents on
opposite sides, less strained, wins: the (\textit{E})-alkene, under thermodynamic control.

\begin{omfigure}
\setchemfig{atom sep=1.9em}
\schemestart
\chemfig{Ph_3P=CH-CH_2CH_3} \+ \chemfig{O=CH-Ph}
\arrow{->}[,0.8]
\chemfig{Ph-[:-30]=[:30]-[:90]CH_2CH_3}
\schemestop

\medskip
\schemestart
\chemfig{Ph_3P=CH-COOEt} \+ \chemfig{O=CH-Ph}
\arrow{->}[,0.8]
\chemfig{Ph-[:-30]=[:30]-[:-30]COOEt}
\schemestop
\par\vspace{6pt}

{\small Two ylides, one aldehyde. Top: the non-stabilised propylidene ylide gives mainly
(\textit{Z})-1-phenylbut-1-ene. Bottom: the \omterm{def:b2:wittig:ylide}{stabilised ylide} gives mainly ethyl
(\textit{E})-3-phenylpropenoate (ethyl cinnamate). Triphenylphosphine oxide, formed in both, is not
shown.}
\end{omfigure}

\section{The Horner--Wadsworth--Emmons reaction}

\begin{definition}[Phosphonates and the HWE reaction]\label{def:b2:wittig:hwe}
A \emph{phosphonate}\index{phosphonate} is an ester of a phosphonic acid, $\mathrm{(RO)_2P({=}O){-}R'}$, with a
\ce{P-C} bond. It is made by the Michaelis--Arbuzov reaction of a trialkyl phosphite with a
halogenoalkane, $\mathrm{P(OEt)_3 + R'Br \to (EtO)_2P(O)R' + EtBr}$. In the
\emph{Horner--Wadsworth--Emmons reaction}\index{Horner--Wadsworth--Emmons reaction} (HWE), the anion of a
phosphonate carrying an electron-withdrawing group on its \ce{CH2} adds to an aldehyde and gives an
alkene and a dialkyl phosphate anion.
\end{definition}

\begin{omfigure}
\setchemfig{atom sep=1.8em}
\schemestart
\chemfig{EtO-P(=[2]O)(-[6]OEt)-CH_2-COOEt}
\arrow{->[\ce{NaH}][$-$ \ce{H2}]}[,0.9]
\chemfig{EtO-P(=[2]O)(-[6]OEt)-\chemabove{C}{\scriptstyle -}H-COOEt}
\schemestop

\medskip
\schemestart
\arrow{->[\ce{PhCHO}]}[,0.9]
\chemfig{*6(=-=(-[:30]=[:-30]-[:30](=[:90]O)-[:-30]O-[:30]Et)-=-)}
\+
\chemfig{EtO-P(=[2]O)(-[6]O^{-})-OEt}
\schemestop
\par\vspace{6pt}

{\small The \omterm{def:b2:wittig:hwe}{Horner--Wadsworth--Emmons reaction} of triethyl phosphonoacetate with benzaldehyde. Sodium
hydride removes the hydrogen between the \omterm{def:b2:wittig:hwe}{phosphonate} and the ester; the anion adds to the aldehyde,
and the four-membered intermediate collapses as in the \omterm{def:b2:wittig:wittig}{Wittig reaction}. The product is ethyl
(\textit{E})-cinnamate; the by-product is the diethyl phosphate anion, soluble in water.}
\end{omfigure}

\begin{proposition}[Why chemists like the HWE reaction]\label{prop:b2:wittig:hwe}
The HWE reaction gives mainly the (\textit{E})-alkene, and its phosphorus by-product, a water-soluble
salt, is removed by a simple aqueous wash; the \omterm{def:b2:wittig:hwe}{phosphonate} anion is also more nucleophilic than the
corresponding \omterm{def:b2:wittig:ylide}{stabilised ylide} and reacts with ketones too.
\end{proposition}

\begin{proof}[Argument]
The \omterm{def:b2:wittig:hwe}{phosphonate} anion is a stabilised carbanion, so its addition is reversible and the reaction runs
under thermodynamic control: (\textit{E}), as with \omterm{def:b2:wittig:ylide}{stabilised ylides}. Unlike a \omterm{def:b2:wittig:ylide}{phosphonium ylide}, it
carries a full negative charge, not compensated by a neighbouring cation: it is more reactive. The
by-product, \ce{(EtO)2PO2-}, is an ionic salt, while triphenylphosphine oxide is a neutral organic
solid that has to be separated from the product by chromatography or crystallisation.
\end{proof}

\section{Choosing a method}

\begin{method}[Wittig disconnection]\label{met:b2:wittig:wittig-disconnection}
\begin{enumerate}
  \item Cut the target \ce{C=C} bond; one carbon becomes a carbonyl carbon, the other the ylide carbon
        (a \omterm{def:b2:wittig:ylide}{phosphonium salt}, from a halide).
  \item Of the two ways of doing it, prefer the one whose halide is primary (bimolecular substitution by
        \ce{Ph3P}) and whose carbonyl compound is an aldehyde.
  \item Choose the reagent from the geometry wanted: a non-stabilised ylide for (\textit{Z}), a
        \omterm{def:b2:wittig:ylide}{stabilised ylide} or a \omterm{def:b2:wittig:hwe}{phosphonate} (HWE) for (\textit{E}).
  \item Check that nothing else in the molecule reacts with the base used.
\end{enumerate}
\end{method}

\begin{method}[Choosing a C=C forming reaction]\label{met:b2:wittig:choose-cc}
\begin{center}
\small
\begin{tabular}{@{}lll@{}}
\hline
Method & Position of the \ce{C=C} & Geometry \\
\hline
E2 or E1 elimination & Zaitsev's rule, mixtures & mostly (\textit{E}), mixtures \\
\omterm{def:b2:enolates-aldol:aldol}{Aldol condensation} & conjugated with \ce{C=O} & (\textit{E}) \\
Wittig, non-stabilised ylide & at the former \ce{C=O} & (\textit{Z}) \\
Wittig, \omterm{def:b2:wittig:ylide}{stabilised ylide}; HWE & at the former \ce{C=O} & (\textit{E}) \\
Alkyne, then poisoned catalyst & at the former triple bond & (\textit{Z}) \\
\hline
\end{tabular}
\end{center}
\medskip
Joining two pieces of similar size at a \ce{C=C}, the \omterm{def:b2:wittig:wittig}{Wittig reaction} and its relatives are usually
the first choice; olefin metathesis, which exchanges the ends of two alkenes, comes in the Year 3
volume.
\end{method}

\begin{history}[The ylide that became a reagent]
Georg Wittig, studying phosphorus ylides in the early 1950s, found that methylidenetriphenylphosphorane
converts benzophenone into 1,1-diphenylethene and triphenylphosphine oxide; with his student Ulrich
Schöllkopf he turned the observation into a general method in 1954. Industry adopted it quickly, for
vitamin A among others, and Wittig shared the 1979 Nobel Prize in Chemistry.
\end{history}

\begin{safety}{\ghs[5mm]{GHS02}\,\ghs[5mm]{GHS05}\,\ghs[5mm]{GHS07}\,\ghs[5mm]{GHS08}}
Butyllithium solutions are flammable and corrosive and react violently with water; sodium hydride
releases hydrogen with water. Triphenylphosphine is harmful and corrosive to the eyes, triethyl
phosphite flammable and harmful; triphenylphosphine oxide is harmful.
% ledger: ghs:BuLi, ghs:PPh3, ghs:triethyl-phosphite, ghs:Ph3PO
\end{safety}

\section{Exercises}

\begin{exercise}[$\star$]\label{exo:b2:wittig:1}
Write the preparation of the ylide \ce{Ph3P=CHCH3} from triphenylphosphine and a halogenoalkane.
\end{exercise}

\begin{exercise}[$\star$]\label{exo:b2:wittig:2}
Give the products of the \omterm{def:b2:wittig:wittig}{Wittig reaction} of benzaldehyde with the ylide made from iodomethane.
\end{exercise}

\begin{exercise}[$\star$]\label{exo:b2:wittig:3}
Classify as stabilised, non-stabilised or aryl-stabilised: \ce{Ph3P=CHCOOEt}, \ce{Ph3P=CHCH3},
\ce{Ph3P=CHCOCH3}, \ce{Ph3P=CHC6H5}. Which needs butyllithium?
\end{exercise}

\begin{exercise}[$\star$]\label{exo:b2:wittig:4}
Give the product of triethyl phosphonoacetate, sodium hydride and propanal, with its geometry.
\end{exercise}

\begin{exercise}[$\star\star$]\label{exo:b2:wittig:5}
Methylenecyclohexane is wanted. Compare a Wittig route from cyclohexanone with the acid-catalysed
dehydration of 1-methylcyclohexan-1-ol.
\end{exercise}

\begin{exercise}[$\star\star$]\label{exo:b2:wittig:6}
Disconnect (\textit{E})-1,2-diphenylethene (stilbene) by the \omterm{def:b2:wittig:wittig}{Wittig reaction}. Why is the Wittig route
a poor choice for pure (\textit{E})-stilbene, and what would you use instead?
\end{exercise}

\begin{exercise}[$\star\star$]\label{exo:b2:wittig:7}
Propanal reacts with \ce{Ph3P=CHCH3} on one side and with \ce{Ph3P=CHCOOEt} on the other. Give the
products and their main geometry.
\end{exercise}

\begin{exercise}[$\star\star$]\label{exo:b2:wittig:8}
Compute the atom economy of the methylenation of cyclohexanone by \ce{Ph3P=CH2}.
\end{exercise}

\begin{exercise}[$\star\star$]\label{exo:b2:wittig:9}
Write the Michaelis--Arbuzov reaction of triethyl phosphite with ethyl bromoethanoate, in two steps:
substitution by phosphorus, then dealkylation by bromide.
\end{exercise}

\begin{exercise}[$\star\star\star$]\label{exo:b2:wittig:10}
Propose a synthesis of 1,6-diphenylhexa-1,3,5-triene from (\textit{E})-but-2-enedial
\ce{OHC-CH=CH-CHO} by two \omterm{def:b2:wittig:wittig}{Wittig reactions}, and say which ylide and how much of it.
\end{exercise}

\begin{exercise}[$\star\star\star$]\label{exo:b2:wittig:11}
(\textit{E})-4-Phenylbut-3-en-2-one can be made by an \omterm{def:b2:enolates-aldol:aldol}{aldol condensation} or by a \omterm{def:b2:wittig:wittig}{Wittig reaction}.
Write both routes and compare them.
\end{exercise}

\begin{exercise}[$\star\star\star$]\label{exo:b2:wittig:12}
Propose two routes to (\textit{Z})-oct-4-ene, one by the \omterm{def:b2:wittig:wittig}{Wittig reaction} and one through an alkyne,
and say which by-products each leaves.
\end{exercise}

\section{Problem: A Pheromone by Wittig}

\begin{problem}[{Weekend problem --- the retrosynthesis of a (Z)-alkene, the non-stabilised ylide and
its selectivity, the cost in triphenylphosphine oxide, and the alkyne alternative}]
\label{pb:b2:wittig:1}
Target: (\textit{Z})-tricos-9-ene, \ce{CH3(CH2)7CH=CH(CH2)12CH3}, the housefly pheromone of the
opening. Molar masses (\unit{g/mol}): C 12.011, H 1.008, O 15.999, P 30.974, Br 79.904.
% ledger: use:muscalure, aw:C, aw:H, aw:O, aw:P, aw:Br, ghs:nonanal, ghs:BuLi

\medskip
\noindent\textbf{Part I --- Retrosynthesis.}
\begin{enumerate}
  \item Write the molecular formula of the target and check that it is an alkene.
  \item Which bond does a Wittig disconnection cut?
  \item Give the two possible pairs (carbonyl compound, \omterm{def:b2:wittig:ylide}{phosphonium salt}).
  \item For the pair nonanal + ylide, which halogenoalkane makes the \omterm{def:b2:wittig:ylide}{phosphonium salt}?
  \item Why is a primary halogenoalkane suitable for this step?
  \item Write the formation of the \omterm{def:b2:wittig:ylide}{phosphonium salt}.
\end{enumerate}

\medskip
\noindent\textbf{Part II --- The ylide and the geometry.}
\begin{enumerate}[resume]
  \item Which base deprotonates this salt? Why would sodium hydroxide not do?
  \item Is the ylide stabilised? What geometry does that predict?
  \item Write the \omterm{def:b2:wittig:wittig}{oxaphosphetane} formed with nonanal.
  \item Why must the reaction be run dry and under inert gas?
  \item Why does the double bond end up exactly between carbons 9 and 10?
  \item What would the acid-catalysed dehydration of tricosan-9-ol give instead?
\end{enumerate}

\medskip
\noindent\textbf{Part III --- Triphenylphosphine oxide.}
\begin{enumerate}[resume]
  \item Write the overall equation from the ylide and nonanal.
  \item Compute the molar masses of the target, of nonanal, of the \omterm{def:b2:wittig:ylide}{phosphonium salt} and of
        triphenylphosphine oxide.
  \item Compute the atom economy of the Wittig step (ylide + aldehyde).
  \item How is triphenylphosphine oxide separated from a hydrocarbon?
  \item Why is the formation of triphenylphosphine oxide the driving force?
  \item Why would the HWE reaction not suit this target?
\end{enumerate}

\medskip
\noindent\textbf{Part IV --- The alkyne route.}
\begin{enumerate}[resume]
  \item Which alkyne, hydrogenated, gives the target, and with which catalyst?
  \item Build that alkyne from dec-1-yne and a halogenoalkane: which halide, which base?
  \item Why must the hydrogenation stop at the alkene, and how does the catalyst ensure it?
  \item Compare the two routes: number of steps, geometry control, by-products.
  \item What is the atom economy of the hydrogenation step?
  \item State the mass of triphenylphosphine oxide co-produced per gram of pheromone by the Wittig
        route.
\end{enumerate}
\end{problem}
